\documentclass[11pt]{article}

\usepackage[margin=1in]{geometry}
\usepackage{amsmath,amssymb,amsthm}
\usepackage{graphicx}
\usepackage[hidelinks]{hyperref}

\newtheorem{theorem}{Theorem}
\newtheorem{lemma}{Lemma}
\newtheorem{remark}{Remark}

\newcommand{\R}{\mathbb R}

\title{Semialgebraic stratified stability for finite measurements}
\author{Agent lane 04}
\date{2026-06-28}

\begin{document}
\maketitle

\begin{abstract}
Alberti, Arroyo and Santacesaria ask whether their low-dimensional inverse
problem theory can be extended to generalized settings with boundary points and
variable parameter number, such as degenerate polygons.  A fixed-dimensional
boundary/corner extension is already recorded in a companion packet, while a
frontier counterexample shows that arbitrary smooth stratified contact is too
broad.  This note gives a positive tame substitute: compact semialgebraic
stratified parameter sets enjoy global H\"older stability from injective
semialgebraic finite measurements.  The proof is an immediate but useful
application of the semialgebraic Lojasiewicz inequality on \(P\times P\).
\end{abstract}

\section{Source question}

The source paper proves H\"older and Lipschitz stability for nonlinear inverse
problems \(F(x)=y\) when the unknown \(x\) lies in a compact subset of a
finite-dimensional differentiable manifold \(M\).  In its conclusion it asks
for extensions beyond manifolds without boundary and fixed dimension.

\begin{center}
\includegraphics[width=.92\linewidth]{figures/open_question_crop.png}
\end{center}

The previous boundary/corner packet handles fixed-dimensional endpoint and
corner charts.  The harder issue is genuine stratification: a positive
dimensional family can collapse to lower-dimensional strata.

\section{Setup}

Let \(P\subset\R^m\) be a compact semialgebraic set.  Think of \(P\) as a
finite-dimensional parameter space that may have strata of different
dimensions.  Let \(X\) be a Banach space and let
\(\iota:P\to X\) be the parameter-to-unknown map.  We assume that \(\iota\) is
\(\alpha\)-H\"older for some \(0<\alpha\le1\):
\[
  \|\iota(p)-\iota(q)\|_X\le L |p-q|^\alpha,\qquad p,q\in P.
\]
Let \(G:P\to\R^N\) be a finite-measurement map.  In applications one may have
\[
  G(p)=QF(\iota(p)),
\]
where \(F\) is the forward map and \(Q\) records finitely many measurements.
The theorem below assumes directly that \(G\) is continuous, semialgebraic, and
injective on \(P\).

\section{Idea of the proof}

The counterexample \(e^{-1/t}\) fails because measurement separation from a
frontier stratum is flatter than every power of parameter separation.
Semialgebraic geometry rules out precisely this phenomenon.  On a compact
semialgebraic set, if a nonnegative semialgebraic function vanishes wherever
another one vanishes, then the first is bounded by a power of the second.  Apply
this to the pairwise parameter distance and pairwise measurement distance on
\(P\times P\).

\section{The theorem}

\begin{theorem}\label{thm:semialgebraic}
Let \(P\subset\R^m\) be compact and semialgebraic, let \(X\) be a Banach space,
and let \(\iota:P\to X\) satisfy
\[
  \|\iota(p)-\iota(q)\|_X\le L |p-q|^\alpha,\qquad p,q\in P,
\]
for some \(L>0\) and \(0<\alpha\le1\).  Let \(G:P\to\R^N\) be continuous,
semialgebraic, and injective.  Then there are constants \(C>0\) and
\(\beta>0\) such that
\[
  \|\iota(p)-\iota(q)\|_X\le C |G(p)-G(q)|^\beta,\qquad p,q\in P.
\]
\end{theorem}

\begin{proof}
Set \(S=P\times P\).  This is a compact semialgebraic subset of
\(\R^{2m}\).  Define
\[
  f(p,q)=|p-q|^2,\qquad g(p,q)=|G(p)-G(q)|^2,\qquad (p,q)\in S.
\]
Both functions are continuous, nonnegative, and semialgebraic.  Since \(G\) is
injective, \(g(p,q)=0\) implies \(p=q\), hence \(f(p,q)=0\).  Thus
\[
  g^{-1}(0)\subset f^{-1}(0).
\]
By the semialgebraic Lojasiewicz inequality on compact sets, there are
constants \(A>0\) and \(\theta>0\) such that
\[
  f(p,q)\le A g(p,q)^\theta,\qquad (p,q)\in S.
\]
Taking square roots gives
\[
  |p-q|\le A^{1/2}|G(p)-G(q)|^\theta.
\]
Therefore
\[
  \|\iota(p)-\iota(q)\|_X
  \le L |p-q|^\alpha
  \le L A^{\alpha/2}|G(p)-G(q)|^{\alpha\theta}.
\]
The result follows with \(C=L A^{\alpha/2}\) and \(\beta=\alpha\theta\).
\end{proof}

\section{Consequences}

\begin{remark}[Derivative injectivity is replaced by tameness plus injectivity]
The theorem does not assume that \(G\) has injective differential on every
stratum.  In the semialgebraic finite-measurement setting, global injectivity is
already enough to force a H\"older inverse on compact parameter sets.  This is
especially useful near singular strata, where the derivative may lose rank even
though the finite measurements still separate all admissible parameters.
\end{remark}

\begin{remark}[Why the frontier counterexample is excluded]
The companion counterexample uses \(F(t)=e^{-1/t}\) near \(t=0\).  This is
\(C^\infty\) on the open stratum and continuous at the frontier, but its
frontier separation is flatter than any power.  Such behaviour is not
semialgebraic.  The present theorem explains why the counterexample is a
genuine smooth-stratified obstruction but not an obstruction in tame algebraic
parameter families.
\end{remark}

\begin{remark}[Relation to degenerate polygons]
If a degenerate polygonal family is encoded by a compact semialgebraic
parameter set \(P\), if the map from parameters to unknowns is H\"older into the
ambient Banach space, and if a finite measurement vector \(G\) is semialgebraic
and injective on \(P\), then the theorem gives the desired H\"older stability.
The hypotheses are intentionally stated at the finite-measurement level; the
paper's arbitrary Banach-valued forward maps need not be semialgebraic in any
finite-dimensional sense.
\end{remark}

\section{Scope and limitations}

This is a partial positive result for the source's generalized-setting
question.  It gives a full H\"older stability theorem for compact
semialgebraic stratified parameter sets with injective semialgebraic finite
measurements.  It does not settle:
\begin{itemize}
\item arbitrary \(C^1\) or \(C^\infty\) stratified frontier degenerations;
\item Banach-valued forward maps without a finite-dimensional tame measurement
  representation;
\item o-minimal structures with exponential functions, where flat behaviour
  like \(e^{-1/t}\) may be definable;
\item the source paper's separate higher-order derivative and nonlinear
  compressed-sensing questions.
\end{itemize}

\section{Novelty check}

The cheap run indexes were searched for arXiv:2009.00574, ``semialgebraic'',
``stratified'', ``Lojasiewicz'', and ``degenerate polygons''.  Existing hits for
this source are the fixed-dimensional boundary/corner packet and the scoped
frontier counterexample.

Web searches on 2026-06-28 for semialgebraic injective maps with H\"older
inverse and Lojasiewicz inequalities found standard semialgebraic metric
results, including Kurdyka--Spodzieja--Szlachcinska, but no separate later
paper applying this directly to the stratified finite-measurement question from
arXiv:2009.00574.

\section{References}

\begin{itemize}
\item G. S. Alberti, A. Arroyo, and M. Santacesaria, \emph{Inverse problems on
  low-dimensional manifolds}, arXiv:2009.00574.
\item K. Kurdyka, S. Spodzieja, and A. Szlachcinska, \emph{Metric properties of
  semialgebraic mappings}, arXiv:1412.5088.
\end{itemize}

\end{document}
